% Propietario conservado: /Users/ruben/Documents/New project/output/EXTREMA_MEDIA_RAZON_RECIPROCIDAD_ES_EN_20260918/ARTICULO/ES/source/libro/reciprocidad_recuperacion.tex
\chapter{Electron coordinate, operator scale, and reversibility of readings}
\label{x_reciprocidad:rec:recuperacion}

\section{Algebraic structure in the complete electronic expression}

The expression in Section~\ref{x_reciprocidad:sec:electron} is now composed with the radial recovery map, retaining its factors and dimensional realization. The electronic prefactor comes from the conditional projectors of addition and multiplication on \(\ell^2(\{1,\ldots,9\}^3)\). On the corresponding principal plane,
\[
 P=\begin{pmatrix}1&0\\0&0\end{pmatrix},\qquad
 Q=\begin{pmatrix}1/4&\sqrt3/4\\\sqrt3/4&3/4\end{pmatrix},
 \qquad\|[P,Q]\|=\frac{\sqrt3}{4}.
\]
The equality follows because
\([P,Q]=(\sqrt3/4)\left(\begin{smallmatrix}0&1\\-1&0\end{smallmatrix}\right)\).
Moreover,
\[
 S=2P-I,\qquad J=\frac4{\sqrt3}[P,Q],\qquad
 S^2=I,\quad J^2=-I,\quad SJ=-JS.
\]
Hence \(n_\pm=(S\pm J)/2\) satisfy \(n_\pm^2=0\). The block retains hyperbolic, elliptic, and nilpotent generators within the algebraic realization already derived. The coefficient \(\sqrt3/4\) and the normalization \(4/\sqrt3\) serve distinct functions.

Using the previously constructed coordinates and regional integers, define
\begin{align}
 \Delta_4&=\frac{\pi-e\log\pi}{270}\left(1+\frac\pi{729}\right),\nonumber\\
 B_e&=\frac{\sqrt3}{4}
       \left[\exp\!\left(\frac\varphi{\pi^2}\right)-22\alpha^3\right]
       (1+15\Delta_4),\qquad
 \kappa_e=80-54\alpha+6\Delta_4.
 \label{x_reciprocidad:rec:eq:electron-basal}
\end{align}
Thus \(\varepsilon_e^{\beta}=B_eR_{\rm act}^{\kappa_e}\). The pre-return action section and its return determine
\[
 \delta_{\rm ret}=H_5-\eta_{\rm ret}=\frac{90}{\pi}\mathcal C_\pi,
 \qquad R_{\rm act}=1+\frac{\delta_{\rm ret}}{\eta_{\rm ret}}.
\]

\begin{corollary}[Radial realization of the electronic coordinate]
\label{x_reciprocidad:rec:cor:electron}
On the branch \(\alpha>0\), \(\eta_{\rm ret}>0\), \(|\xi|<1\),
\begin{align}
 \varepsilon_e^{\beta}={}&
 \frac{\sqrt3}{4}
 \left[\exp\!\left(\frac\varphi{\pi^2}\right)-22\alpha^3\right]
 (1+15\Delta_4)\nonumber\\
 &\times\left[1+\frac{90\mathcal C_\pi}{\pi\alpha}
 \frac{1+R_\star^4}{499-501R_\star^4}\right]^{80-54\alpha+6\Delta_4}.
 \label{x_reciprocidad:rec:eq:electron-radio}
\end{align}
In the Klein coordinate of the marked chart,
\begin{equation}
 R_{\rm act}=1+
 \frac{\delta_{\rm ret}}{\alpha(-500x_K-1)}.
 \label{x_reciprocidad:rec:eq:accion-klein}
\end{equation}
The corresponding domains are
\[
 \frac1{500}<\xi<1,\qquad
 0<R_\star^4<\frac{499}{501},\qquad
 -1<x_K<-\frac1{500}.
\]
\end{corollary}
\begin{proof}
We substitute \eqref{x_reciprocidad:rec:eq:eta-radio} into \(R_{\rm act}=1+\delta_{\rm ret}/\eta_{\rm ret}\) and then into the complete electronic expression \eqref{x_reciprocidad:eq:el-formula-completa}. The Klein chart gives \(\xi=-x_K\). The intervals are obtained from \(\eta_{\rm ret}=\alpha(500\xi-1)>0\) and the already proved monotone correspondences.
\end{proof}

The substitution preserves the regional deficit, torsional normalization, return character, and incompatibility norm. It expresses the same coordinate through a recoverable geometric reader. The radius and \(\alpha\) remain correlated by their generation; they are not independent parameters fitted to a mass. The dimensional realization remains
\[
 E_e^{\beta}=\varepsilon_e^{\beta}\mathcal U_E,\qquad
 m_e^{\beta}=E_e^{\beta}/c_{\rm int}^2.
\]
The action decade already incorporated into the Planck sections is not multiplied again onto a coordinate expressed in the energy chart.

A second exact form, useful for isolating the correction, is
\[
 \widehat\delta_{\rm ret}=\frac{90\mathcal C_\pi}{\pi H_5},\qquad
 R_{\rm act}=(1-\widehat\delta_{\rm ret})^{-1},\qquad
 \varepsilon_e^{\beta}=B_e(1-\widehat\delta_{\rm ret})^{-\kappa_e}.
\]
Positivity of both sections ensures \(1-\widehat\delta_{\rm ret}>0\). The three expressions are identities of the same return, not independent predictions.

\section{Separation of scale and shape in a positive mass fibre}

The operator law of return receives a finite positive fibre of routes together with its readers. To formulate its matrix consequence independently of species names, let \(\mathcal H\) have dimension \(d\ge1\), let \(B=B^*\) be the exponent operator, and let \(M^0\) be positive definite. The return is the congruence
\begin{equation}
 M^\beta=R_{\rm act}^{B/2}M^0R_{\rm act}^{B/2}.
 \label{x_reciprocidad:rec:eq:masa-operador}
\end{equation}
In the route basis of the diagonal reader, \(Be_\gamma=\kappa_\gamma e_\gamma\) and
\(M^0e_\gamma=m_e^{(0)}\chi_{\rm ref}(\sigma_\gamma-\sigma_e,\eta_\gamma-\eta_e)e_\gamma\). The signatures and character are those of the route construction; the following theorem applies to each positive fibre without replacing its selector.

\begin{theorem}[Uniform and unimodular decomposition of return]
\label{x_reciprocidad:rec:thm:masa}
If \(R_{\rm act}>0\), define
\[
 \bar\kappa=\frac{\tr B}{d},\qquad B_0=B-\bar\kappa I,
 \qquad S_0=R_{\rm act}^{B_0/2}.
\]
Then
\begin{align}
 M^\beta&=R_{\rm act}^{\bar\kappa}S_0M^0S_0,
 \qquad\det S_0=1,\label{x_reciprocidad:rec:eq:masa-forma}\\
 \det M^\beta&=R_{\rm act}^{\tr B}\det M^0.\label{x_reciprocidad:rec:eq:masa-det}
\end{align}
With \(\widehat M=M/(\det M)^{1/d}\),
\begin{equation}
 \widehat M^\beta=S_0\widehat M^0S_0.
 \label{x_reciprocidad:rec:eq:masa-normalizada}
\end{equation}
\end{theorem}
\begin{proof}
The scalar part of \(B\) commutes with \(B_0\). Functional calculus gives
\(R_{\rm act}^{B/2}=R_{\rm act}^{\bar\kappa/2}S_0\). In an orthonormal eigenbasis of \(B\),
\[
 \det S_0=R_{\rm act}^{\frac12\sum_\gamma(\kappa_\gamma-\bar\kappa)}=1.
\]
Multiplication of determinants proves \eqref{x_reciprocidad:rec:eq:masa-det}. Dividing \eqref{x_reciprocidad:rec:eq:masa-forma} by the \(d\)-th root of that determinant cancels the scalar factor. The proof does not require \(B\) to commute with \(M^0\); it does require positive definiteness for normalization by the determinant.
\end{proof}

In the joint diagonal basis of the source construction, the individual ratios are
\[
 q_\gamma=\frac{m_\gamma^\beta}{m_\gamma^{(0)}}R_{\rm act}^{-\bar\kappa}
 =R_{\rm act}^{\kappa_\gamma-\bar\kappa},\qquad
 \prod_\gamma q_\gamma=1.
\]
The formal inversion \(R_{\rm act}\mapsto R_{\rm act}^{-1}\) inverts each \(q_\gamma\). It differs from oriented radial reciprocity, which preserves \(R_{\rm act}\) by \eqref{x_reciprocidad:rec:eq:eta-orientada}. The congruence preserves the normalized determinant, not the individual eigenvalues or the sum of masses. The null sector lies outside this positive normalization. For \(d=1\), \(B_0=0\): the electron's return lies entirely in the uniform factor and remains in \eqref{x_reciprocidad:rec:eq:electron-radio}.

Substitution of \eqref{x_reciprocidad:rec:eq:accion-klein} also acts on \eqref{x_reciprocidad:rec:eq:masa-operador} through functional calculus. If
\[
 q(\xi)=1+\frac{\delta_{\rm ret}}{\alpha(500\xi-1)}>0,
\]
then
\[
 M^\beta=e^{\frac12(\log q(\xi))B}\,M^0\,
          e^{\frac12(\log q(\xi))B}.
\]
The genealogy of each exponent and each multisection is preserved. The structure shared with the flow of \(K\) consists of removing a uniform component and then applying a deformation of determinant one; \(B_0\) is not identified with \(\operatorname{diag}(u)\), nor \(R_{\rm act}\) with \(\rho\).

\section{Circular compression and identification of transport}

\begin{theorem}[Bilateral recovery of a unitary operator]
\label{x_reciprocidad:rec:thm:operador}
Let \(C\) be unitary on a Hilbert space, and let
\[
 T=\frac{8I+C}{9},\qquad D=\frac{\sqrt8(C-I)}9.
\]
Then
\begin{equation}
 T^*T+D^*D=I,\qquad
 v=Tv-\frac{Dv}{\sqrt8},\qquad Cv=Tv+\sqrt8Dv.
 \label{x_reciprocidad:rec:eq:unitario}
\end{equation}
\end{theorem}
\begin{proof}
Expansion gives
\[
 81T^*T=65I+8(C+C^*),\qquad
 81D^*D=16I-8(C+C^*).
\]
The sum is \(81I\). The other identities follow by taking linear combinations of the two definitions. At a circular eigenvalue \(|z|=1\), the balance is
\(1-|(8+z)/9|^2=(8/81)|1-z|^2\).
\end{proof}

The archived pair recovers both the state and its image. For the invertible frame \(O_K=(K,SK,\ldots,S^{11}K)\) already constructed, the twelve responses form
\[
 Y=(CK,CSK,\ldots,CS^{11}K),\qquad C=YO_K^{-1}.
\]
If \(CS=SC\), a single vector response \(y=CK\) determines all the columns: \(Y=(y,Sy,\ldots,S^{11}y)\). The general case requires twelve vector trials. Recovery does not reduce that information to a single scalar.

The real diagonal generator \(H_K=\operatorname{diag}(h_i)\) admits the unitary chart
\[
 C_K=(H_K+iI/2)(H_K-iI/2)^{-1}.
\]
Each real eigenvalue is transformed into a number of modulus one. Moreover,
\[
 C_K-I=i(H_K-iI/2)^{-1},\qquad
 H_K=\frac i2(C_K+I)(C_K-I)^{-1}.
\]
Both inverses exist because \(H_K\) is self-adjoint and its spectrum is real. The archive of \eqref{x_reciprocidad:rec:eq:unitario} therefore recovers this finite realization of the direction of \(K\). Its dimension twelve and its domain remain distinct from those of a spectral operator for zeta.

\section{Hierarchy of invariants and scope of the proofs}

The completion \(1000=729+271\), with normalized reading \(1=729/1000+271/1000\), retains its generative function. The subsequent conservation laws are organized according to the object transported:
\begin{center}
\small
\begin{tabularx}{\linewidth}{@{}p{.27\linewidth}Xp{.25\linewidth}@{}}
\toprule
Realization & Invariant & Recovery or reciprocity\\
\midrule
Archivo bilateral & Joint norm of reading and memory & State and image under the operator\\
Action ellipse & Product \(a_Sb_S=2\hbar\) & Radius and angular orientation\\
Flow of \(K\) & Product of twelve radii & Inversion \(t\mapsto-t\)\\
Deformed marked lattice & Covolume and ternary isometry & \(\Lambda_t^*=\Lambda_{-t}\)\\
Theta--Mellin reading & Poles and residues & Entire difference and functional equation\\
Positive mass fibre & Determinant after removing the common scale & Congruence by \(S_0\)\\
\bottomrule
\end{tabularx}
\end{center}

Conservation of volume does not imply conservation of each vector's norm. Polar conservation permits variation in the entire part. Recovery of action requires orientation to be retained. These distinctions express the operational content of evolutionary conservation: each transformation retains certain data, changes others, and preserves the register that allows the correspondence to be reconstructed.

The composition extends the function of \(K\) already proved in the article. The register participates in the reading of \(\alpha\); its direction is recovered from memory; that direction determines a radial geodesic and a reciprocal lattice deformation; the lattice family has a stable spectral reading. The action section provides the link to the electronic coordinate and to the positive congruence of the mass operator. This yields a chain of maps with explicit domains, inverses and invariants.

The supplement distinguishes written proofs from formalized ones. \texttt{FlujoReciproco.lean} contains ten theorems on diagonal flow, centering, determinants, inversion and lifting to twenty-four dimensions. \texttt{RadioKlein.lean} contains eleven rational theorems on charts, bounds, reciprocities and oriented recovery. The hyperbolic proofs, complete substitutions, separation of the mass operator and theta--Mellin argument appear in the text; their formalization is not attributed to compilation of those two files. Documentary receipts record provenance and reproduction, with a scope distinct from that of a mathematical proof.
